\scanpage{35}

% B01 | top centre
\mspage{\unsure{5|S}}

% B02 | upper right
GRÖTHENDIECK.

% B03 | printed letterhead, upper left
% letterhead
\begin{center}
UNIVERSITÉ DE PARIS\\
---$\cdot$---\\
ÉCOLE NORMALE SUPÉRIEURE\\
---$\cdot$---\\
45, RUE D'ULM\\
ODÉON 06$\cdot$45\\
---$\cdot$---
\end{center}

% B04 | theorem, left label and seven lines
\uuline{Th 1} Soit $E$ un ensemble simplicial, \deleted{\illegible{1}} $G$ un groupe discret qui opère sur $E$ \deleted{\illegible{1}}par des transformations \unsure{simplicials|simpliciales}, $B\inserted{=E/G}$ l'ensemble simplicial quotient, \deleted{\illegible{1}} $p:E\to B$ la projection canonique. Si $E$ est connexe, l'image de $p_*:\pi_1(E)\longrightarrow\pi_1(B)$ est un sousgroupe distingué de $\pi_1(B)$, et $\pi_1(B)/\operatorname{im}p_*(\pi_1(E))=G/S$, où $S$ est le sousgroupe (distingué) de $G$ engendré par les stabilisateurs des sommets de $E$, et $E/S$ est le revêtement \deleted{\unsure{universel|galoisien}} de $B$ correspondant à \inserted{l'image de $p_*(\pi_1(E))$}

% B05 | Déf., two lines under the theorem
\uuline{Déf.} une arête de $E$ est un 1-simplexe de $E$ muni d'un signe $+$ ou $-$. Si $\sigma$ est un 1-simplexe, on pose $\partial_1(+\sigma)=\partial_1(\sigma)$, $\partial_0(+\sigma)=\partial_0(\sigma)$, $\partial_1(-\sigma)=\partial_0(\sigma)$, $\partial_{\unsure{0|d}}(-\sigma)=\partial_1(\sigma)$.

% B06 | lower paragraph, Un chemin through suivants
Un chemin est une suite \inserted{d'arêtes} $\gamma=(\gamma_1,\ldots,\gamma_n)$ telle que $\partial_1(\gamma_{i+1})=\partial_0(\gamma_i)$. On pose $\partial_1(\gamma)=\partial_1(\gamma_1)$ et $\partial_0(\gamma)=\partial_0(\gamma_n)$. Un \deleted{point} \inserted{sommet} de base $x_0$ étant fixé dans $E$, un lacet est un chemin $\gamma$ tel que $\partial_1(\gamma)=\partial_0(\gamma)=x_0$. 2 chemins $\gamma$ et $\gamma'$ sont homotopes si leurs réalisations géométriques sont homotopes modulo leurs extrémités. Cette relation d'équivalence est compatible avec la loi de composition « mise bout à bout » des chemins, et est engendrée par les suivants :

% archive sheet number: 34

